% !TeX root = ../thesis.tex
% Kapitel 5
\chapter{Experimentelle Evaluation}
\label{cha:evaluation}

\section{Versuchsumgebung}
\label{sec:setting}
\todo{Parameter aus \code{experiment.yaml} und verwendeter Lauf.}

\section{Experimente}
\label{sec:experimente}
\todo{Experimente den Teilfragen zuordnen, bevor die Ergebnisse vorliegen.}

\section{Ergebnisse}
\label{sec:ergebnisse}
\begin{table}[htbp]
  \centering
  \caption[Vergleich der Methoden]{Methoden und Selektoren bei Maskierungsrate 0
    (Mittelwert $\pm$ Standardabweichung über 5 Durchgänge, Platzhalter).}
  \label{tab:ergebnisse}
  \begin{tabular}{llccc}
    \toprule
    Methode & Selektor & P@3 & R@3 & ECE \\
    \midrule
    baseline   & none   & -- & -- & -- \\
    knn\_gower & none   & -- & -- & -- \\
    knn\_gower & grande & -- & -- & -- \\
    clustering & none   & -- & -- & -- \\
    clustering & grande & -- & -- & -- \\
    \bottomrule
  \end{tabular}
\end{table}
\todo{Ergebnisse eintragen, Verlauf über Maskierungsraten, Signifikanz (Friedman und Wilcoxon).}

\section{Diskussion}
\label{sec:diskussion}
\todo{Teilfragen beantworten, Validität und Grenzen.}
